
@@M0@@ sur @@M1@@; notons @@M2@@ l'homomorphisme défini par la suite @@M3@@, de sorte que @@M4@@ est déduit de @@M5@@ par changement de base de @@M6@@ à @@M7@@. Comme la suite @@M8@@ est régulière (au sens classique), l'homomorphisme @@M9@@ est régulier, et on est sous les conditions du lemme, appliqué au changement de base @@M10@@. Or les @@M11@@ sont libres de type fini sur @@M12@@, donc on peut appliquer les résultats du a). Par suite @@M13@@, et @@M14@@ est libre de type fini.

\textbf{Proposition 1.3.} \uline{Soient} @@M15@@ \uline{un anneau}, @@M16@@ \uline{un} @@M17@@-\uline{module projectif de type fini}, @@M18@@ \uline{un homomorphisme régulier}, @@M19@@. \uline{Alors l'idéal} @@M20@@ \uline{est quasi-régulier} (EGA IV 16.9.7), \uline{c'est-à-dire} (EGA IV 16.9.8):
\begin{enumerate}[label=\roman*)]
\item @@M21@@ \uline{est de type fini}.
\item @@M22@@ \uline{est un} @@M23@@-\uline{module projectif de type fini};
\item \uline{l'homomorphisme canonique}
\end{enumerate}
@@M24@@
\uline{est un isomorphisme}.

On a vu dans la démonstration de 1.2 b), que @@M25@@ est projectif de type fini. Reprenons les notations de 1.2 b) (la question étant locale sur @@M26@@, on suppose encore @@M27@@ libre de type fini). Alors l'homomorphisme canonique (1.3.1) se déduit de l'homomorphisme analogue pour @@M28@@ par le changement de base @@M29@@, d'après les relations 1.2. Or ce dernier est un isomorphisme, donc il en est de même pour (1.3.1).

\textbf{Définition 1.4.} \uline{Soient} @@M30@@ \uline{un site commutativement localement annelé}, @@M31@@ \uline{un} @@M32@@-\uline{Module localement libre de type fini}, @@M33@@ \uline{un homomorphisme} @@M34@@-\uline{linéaire}. \uline{On dit que} @@M35@@ \uline{est} régulier \uline{si le complexe de Koszul} @@M36@@ \uline{défini par} @@M37@@ \uline{est acyclique en degrés} @@M38@@. \uline{Soit} @@M39@@ \uline{un Idéal de} @@M40@@; \uline{on dit que} @@M41@@ \uline{est}
